%!TEX root = ../../thesis.tex
% Generated by scripts/evaluate.py -- do not edit; edit the emitter instead.
\begin{figure}[tbp]
  \centering
  \begin{subfigure}[b]{0.99\linewidth}
    \centering
    \begin{tikzpicture}
      \begin{axis}[thesis result plot, width=0.95\linewidth, height=4.2cm,
        xlabel={learner iteration}, ylabel={structural byte rate},
        xmin=0, ymin=0, ymax=1]
      \addplot [KIT line plot A, forget plot]
        table [col sep=comma, x=iteration, y=s0] {figures/data/structural-rate-cjson.csv};
      \addplot [KIT line plot B, forget plot]
        table [col sep=comma, x=iteration, y=s1] {figures/data/structural-rate-cjson.csv};
      \addplot [KIT line plot C, forget plot]
        table [col sep=comma, x=iteration, y=s2] {figures/data/structural-rate-cjson.csv};
      \addplot [KIT line plot D, forget plot]
        table [col sep=comma, x=iteration, y=s3] {figures/data/structural-rate-cjson.csv};
      \addplot [KITblack50, dashed, forget plot] coordinates {(0,0.0860) (191,0.0860)};
      \end{axis}
    \end{tikzpicture}
    \caption{cJSON}
  \end{subfigure}
  \\[1.4ex]
  \begin{subfigure}[b]{0.99\linewidth}
    \centering
    \begin{tikzpicture}
      \begin{axis}[thesis result plot, width=0.95\linewidth, height=4.2cm,
        xlabel={learner iteration}, ylabel={structural byte rate},
        xmin=0, ymin=0, ymax=1]
      \addplot [KIT line plot A, forget plot]
        table [col sep=comma, x=iteration, y=s0] {figures/data/structural-rate-libxml2.csv};
      \addplot [KIT line plot B, forget plot]
        table [col sep=comma, x=iteration, y=s1] {figures/data/structural-rate-libxml2.csv};
      \addplot [KIT line plot C, forget plot]
        table [col sep=comma, x=iteration, y=s2] {figures/data/structural-rate-libxml2.csv};
      \addplot [KIT line plot D, forget plot]
        table [col sep=comma, x=iteration, y=s3] {figures/data/structural-rate-libxml2.csv};
      \addplot [KITblack50, dashed, forget plot] coordinates {(0,0.0544) (399,0.0544)};
      \end{axis}
    \end{tikzpicture}
    \caption{libxml2}
  \end{subfigure}
  \\[1.4ex]
  \begin{subfigure}[b]{0.99\linewidth}
    \centering
    \begin{tikzpicture}
      \begin{axis}[thesis result plot, width=0.95\linewidth, height=4.2cm,
        xlabel={learner iteration}, ylabel={structural byte rate},
        xmin=0, ymin=0, ymax=1]
      \addplot [KIT line plot A, forget plot]
        table [col sep=comma, x=iteration, y=s0] {figures/data/structural-rate-picohttpparser.csv};
      \addplot [KIT line plot B, forget plot]
        table [col sep=comma, x=iteration, y=s1] {figures/data/structural-rate-picohttpparser.csv};
      \addplot [KIT line plot C, forget plot]
        table [col sep=comma, x=iteration, y=s2] {figures/data/structural-rate-picohttpparser.csv};
      \addplot [KIT line plot D, forget plot]
        table [col sep=comma, x=iteration, y=s3] {figures/data/structural-rate-picohttpparser.csv};
      \addplot [KITblack50, dashed, forget plot] coordinates {(0,0.0281) (399,0.0281)};
      \end{axis}
    \end{tikzpicture}
    \caption{picohttpparser}
  \end{subfigure}
  \caption[Structural byte rate against training iteration]{
    Fraction of self-play actions drawn from the target's grammar-relevant byte
    set, one line per seed.
    The dashed line is the untrained rate measured at iteration 0.
    All three targets rise far above it, and all three peak early and then decay~(\cref{fig:structural-decay}).}
  \label{fig:structural-rate}
\end{figure}
